% 探针6：二分定位 underfull 幻影断行来源（裸段 / multicols / ansblock 三变体）
\def\mthreepure{0}
\documentclass[fontset=none]{ctexart}
\usepackage{qp-m3}
\usepackage{pifont}%
\xeCJKDeclareCharClass{Default}{"2212}%
\xeCJKDeclareCharClass{Default}{"2013}%
\xeCJKsetup{CJKglue={\hskip 0.10em plus 0.08em minus 0.05em}}%
\begin{document}
% v1 裸段
故 \(|OA|\cdot|OB|=\sqrt{x_{A}x_{B}}\cdot\sqrt{x_{A}x_{B}+2p(x_{A}+x_{B})+4p^{2}}\)，代入韦达读数得 \((p/2)\cdot\sqrt{p^{2}/4+6p^{2}+4p^{2}}=(p/2)\cdot(\sqrt{41}/2)p=\sqrt{41}p^{2}/4\)．

\begin{multicols}{2}
% v2 multicols 内裸段
又 \(|CN|=\sqrt{(x_{0}-5)^{2}+y_{0}^{2}}=\sqrt{(x_{0}-5)^{2}+8x_{0}}\)，即 \(\sqrt{(x_{0}-1)^{2}+24}\geqslant 2\sqrt{6}\)，当且仅当 \(x_{0}=1\)（\(N(1,\pm 2\sqrt{2})\)）且 \(C\)、\(M\)、\(N\) 三点共线（\(M\) 为线段 \(CN\) 与圆的交点）时取等——取等可行：\(C(5,0)\) 在圆内（\(|CB|=1<2\)）而 \(N\) 在圆外（\(|NB|^{2}=(x_{0}-6)^{2}+8x_{0}=(x_{0}-2)^{2}+32>4\) 恒成立），线段 \(CN\) 必与圆相交．最小值 \(2\sqrt{6}\)，选 D．
{\ansblockgrayfalse
\begin{ansblock}[probe-v3]
% v3 ansblock（括线）内同段
又 \(|CN|=\sqrt{(x_{0}-5)^{2}+y_{0}^{2}}=\sqrt{(x_{0}-5)^{2}+8x_{0}}\)，即 \(\sqrt{(x_{0}-1)^{2}+24}\geqslant 2\sqrt{6}\)，当且仅当 \(x_{0}=1\)（\(N(1,\pm 2\sqrt{2})\)）且 \(C\)、\(M\)、\(N\) 三点共线（\(M\) 为线段 \(CN\) 与圆的交点）时取等——取等可行：\(C(5,0)\) 在圆内（\(|CB|=1<2\)）而 \(N\) 在圆外（\(|NB|^{2}=(x_{0}-6)^{2}+8x_{0}=(x_{0}-2)^{2}+32>4\) 恒成立），线段 \(CN\) 必与圆相交．最小值 \(2\sqrt{6}\)，选 D．
\end{ansblock}}
\end{multicols}
\end{document}
